\section{Augmentation Techniques Backgrounds}\label{Augmentations}

This section introduces four groups of commonly used image augmentation techniques in computer vision: geometric, color-space, cutout and mixed. Each technique is demonstrated using three images of birds, including Genus Fulmarus (Eissturmvogel), Common Goldeneye (Schnellente), and Great Cormorant (Kormoran). These three images are selected to demonstrate the effects in scenarios where the object of interest is clear (Kormoran), the object shape is small (Eissturmvogel), and when working with a group of birds (Schnellente).
\subsection{Geometric augmentation}
Geometric augmentations are a type of image augmentation technique used in computer vision that modify the geometric properties of an image, such as its size, position, orientation, and perspective. In this section, we will focus on five common geometric augmentations that are particularly useful for aerial images: Cropout, Flipping, Motion Blur, Perspective and Rotation. 
\subsubsection{CenterCrop} \label{CenterCrop}
\begin{figure}
	 \centerline{\includegraphics[width=1\textwidth]{src/augmentations/CenterCrop.png}}
	 {\caption{Example of CenterCrop with proportion of 0.2 and 0.4 }\label{fig:CenterCrop}}
\end{figure}

CenterCrop transformation is a type of image augmentation technique that involves cropping an image from its center to generate a smaller sub-image. This transformation can be used to remove irrelevant areas from the image and focus on the most important features. The CenterCrop transformation can also help to standardize the size of images used for training machine learning models. The CenterCrop transformation can be advantageous for aerial images, as it can eliminate irrelevant background areas, enabling the focus on essential features. This is particularly useful for our dataset, which contains numerous images with center biases. 

However, it is crucial to note that in some pictures where there is no center bias, such as the Kormoran image in Figure \ref{fig:CenterCrop} or images showing a flock of birds, this transformation may remove important information from the image. Therefore, it is necessary to evaluate each image carefully and determine whether the CenterCrop transformation is suitable, as its usage may depend on the image's characteristics.

In our experiment, we will apply CenterCrop transformation to our aerial images, using proportions of 0.2 and 0.4, as shown in Figure \ref{fig:CenterCrop}. This means that the center of the image will be cropped by one-fifth and two-fifths of its width and height. Then the image will be resized to the original size. 


\subsubsection{Flipping} \label{Flipping}
\begin{figure}
	 \centerline{\includegraphics[width=0.8\textwidth]{src/augmentations/flipping.png}}
	 {\caption{Examples of Flipping}\label{fig:flipping}}
\end{figure}
Flipping transformation is an image augmentation technique that flips an image horizontally or vertically. It involves reflecting the image along a vertical or horizontal axis, which results in a mirror image of the original image. 

The advantage of flipping is that it is a simple transformation that does not risk any loss of information. However, it is limited in terms of variety because it can only create mirrored versions of existing examples, which may not be sufficient for capturing the full range of variations in the dataset. Figure \ref{fig:flipping} illustrates the efffect of the transformations.

\subsubsection{Motion Blur} \label{MotionBlur}
Motion Blur transformation is an image augmentation technique that simulates Motion Blur in an image. The Motion Blur transformation works by applying a blur filter to the image, which simulates the effect of an object in motion. This is accomplished by convolving the image with a kernel function that creates the blur effect. Figure \ref{fig:motion_blur} illustrated the blur transformation. The size and direction of the kernel function can be varied to control the strength and direction of the Motion Blur effect. The blur limit in Motion Blur augmentation refers to the maximum amount of blur that can be applied to an image. This limit can be defined in terms of the size and direction of the kernel function used to apply the blur effect. 

\begin{figure}[!h]
	 \centerline{\includegraphics[width=0.8\textwidth]{src/augmentations/motion_blur.png}}
	 {\caption{Examples of Motion Blur with blur litmit of 30} \label{fig:motion_blur}}
\end{figure}


One advantage of Motion Blur transformation is that it can help to generate a more diverse set of training images for machine learning models. By varying the size and direction of the kernel function, it is possible to generate a large number of images with different types and strengths of Motion Blur. This can help improve the accuracy and generalization of the trained model.

While Motion Blur transformation can be useful for generating a diverse set of training images, it does have some significant drawbacks, particularly for object detection in aerial images of marine wildlife. A significant disadvantage of this technique is that it can reduce the image's sharpness and detail, making it more challenging to detect or recognize objects in the image. In the case of aerial images of marine wildlife, where shape and detail are essential for accurate detection, Motion Blur transformation can negatively impact the performance of the machine learning model. 

\subsubsection{Perspective} \label{Perspective}

Perspective transformation is an image augmentation technique that involves applying a geometric transformation to an image in order to alter its perspective. To perform perspective transformation, the algorithm first defines a set of four source points and a set of four destination points that define the corners of the original and transformed image, respectively. Then, it computes a transformation matrix that maps the source points to the destination points using techniques such as homography or projective geometry. An example of perspective transformation is the transformation of a squared image of a square into a trapezoid shape, then resize it to their original size. This transformation can create the illusion of a 3D perspective in a 2D image. 

Hyperparameter of Perspective includes random scale factor. The scale factor is used to determine the new points' distance from the original image corners during the transformation, which is calculated by sampling from a normal distribution with a mean of zero and a standard deviation set by the scale factor. A larger scale factor leads to a more pronounced transformation.  In the experiment, the scale factor will be tuned, the effect of Perspective transformation is shown in Figure \ref{fig:perspective}. 

\begin{figure}[!h]
	 \centerline{\includegraphics[width=0.8\textwidth]{src/augmentations/Perspective Transformation.png}}
	 {\caption{Example of Perspective with random scale between 0.1 and 0.2}\label{fig:perspective}}
\end{figure}

It should be noted that in the case of aerial images, perspective transformation may not be as effective. Aerial images are typically captured from a top-down perspective, where the camera is fixed in a position and oriented straight downward. As a result, the scene in the image is already in a consistent perspective, and applying a perspective transformation may not yield significant benefits.

\subsubsection{Rotation} \label{Rotation}

Rotation transformation is an image augmentation technique that involves rotating an image by a certain angle. The rotation can be performed in either a clockwise or counter-clockwise direction.


One advantage of rotation transformation is that it can increase the diversity of the training data without risking any information loss. This is particular useful for our dataset because marine wildlife aerial images often have objects in different orientations and angles. By applying rotation transformations, the model can learn to recognize the same object in various angles, which improves the model's robustness and accuracy in real-world scenarios. Additionally, rotation transformation is a simple and computationally efficient technique that can generate a wide variety of training examples. In this thesis, we apply a rotation transformation to our images within a range of -90 to 90 degrees. The impact of this transformation is demonstrated in Figure \ref{fig:rotation}.

\begin{figure}[!h]
	 \centerline{\includegraphics[width=0.8\textwidth]{src/augmentations/rotation.png}}
	 {\caption{Examples of Rotation}\label{fig:rotation}}
\end{figure}
\subsection{Color-space augmentation}
Color-space augmentation is an image augmentation technique that involves altering the color channels of an image to create new training images for machine learning models. This technique can be used to modify the color distribution of images, introduce color variations, and aid models in adapting to various lighting conditions. In this section, we will present Contrast Limited Adaptive Histogram Equalization (CLAHE), Sharpen, Channel Shuffling, Color Jitter, Solarize, ToGray, ToSepia, and Gaussian Noise Injection as examples of color-space augmentation.
\subsubsection{CLAHE} \label{CLAHE}

Contrast Limited Adaptive Histogram Equalization (CLAHE) is an image enhancement technique that improves the contrast and brightness of an image. This technique is widely used in image processing and computer vision tasks to improve the visual quality of images and make it easier to detect objects and features, example of CLAHE is shown in Figure \ref{fig:CLAHE}.

\begin{figure}[!h]
	 \centerline{\includegraphics[width=0.8\textwidth]{src/augmentations/CLAEH.png}}
	 {\caption{Examples of CLAHE with contrast limit between 1 and 4 and title grid with size of 8} \label{fig:CLAHE}}
\end{figure}

CLAHE works by dividing the image into small rectangular regions called tiles and then equalizing the histogram of each tile separately. This approach ensures that the brightness and contrast of each tile are optimized independently, which can be beneficial in cases where the brightness or contrast varies significantly across the image.

One advantage of CLAHE is that it is an adaptive technique, meaning that it can adapt to the local characteristics of the image. This allows it to provide better results than traditional histogram equalization, which applies the same transformation to the entire image. CLAHE also allows for better visualization of fine details and edges in images, which can be especially useful in medical imaging applications or aerial images of marine wildlife where subtle changes in texture or contrast can be important for detecting features or anomalies.

\subsubsection{Sharpen} \label{Sharpen}

Sharpen augmentation is a technique used in image processing and computer vision that enhances the edges and fine details in an image. The technique involves applying a filter to the image that increases the contrast between adjacent pixels, making the edges more distinct and prominent.

Sharpen augmentation can be particularly advantageous in the context of aerial images for a number of reasons. Aerial images of the marine environment, for instance, often contain low-contrast regions such as large bodies of water or sky that can reduce the sharpness and clarity of the image. Sharpen augmentation can help to increase the contrast between adjacent pixels, making edges and fine details more visible and distinct. This can improve the visual quality of the image and make it easier to detect and recognize objects, including marine wildlife, in the image. 

The alpha parameter in sharpen augmentation is a factor that controls the strength of the filter. A higher value of alpha will result in a stronger sharpening effect, while a lower value will result in a weaker effect.  The lightness parameter in sharpen augmentation is a factor that controls the amount of lightness or brightness in the image. A higher value of lightness will increase the brightness of the image, while a lower value will decrease it. Examples of augmented images are represented in Figure \ref{fig:Sharpen}

\begin{figure}[!h]
	 \centerline{\includegraphics[width=0.8\textwidth]{src/augmentations/Sharpen.png}}
	 {\caption{Examples of Sharpen with alpha value between 0.2 and 0.5, lightness value between 0.5 and 1.0} \label{fig:Sharpen}}
\end{figure}

\subsubsection{Channel Shuffling}\label{ChannelShuffles}
Channel Shuffling augmentation is an image transformation technique that can be used to generate a diverse set of training images for machine learning models. In channel shuffles, the color channels of an image are randomly permuted, which can introduce a significant amount of variation to the image data.

The advantage of Channel shuffling is its simplicity in both understanding and implementation. However, the variation of the transformation is limited to the permutation of the channels, which is only six. This limitation restricts the ability to enhance the robustness of the data. Moreover, one color channel may be more important than others. For example, if blue is an important feature to recognize the ocean, shuffling the channels could pose an obstacle for the model. Figure \ref{fig:ChannelShuffles} show the effect of Channel Shuffling on the image.
\begin{figure}[!h]
	 \centerline{\includegraphics[width=0.8\textwidth]{src/augmentations/ChannelShuffle.png}}
	 {\caption{Examples of Channel Shuffling} \label{fig:ChannelShuffles}}
\end{figure}

\subsubsection{Color Jitter}\label{ColorJitter}
Color Jitter is an image augmentation technique that introduces random color variation to an image by perturbing the color channels of the image. The goal of this technique is to increase the diversity of the training data-set for machine learning models, which can help improve the model's ability to generalize to new, unseen data.

Color Jitter works by randomly perturbing the hue, saturation, brightness, and contrast of an image. These perturbations are generated by adding small amounts of noise or by scaling the color values in channels. By applying these perturbations, Color Jitter creates new images that are visually distinct from the original image, which can help improve the model's ability to detect and recognize objects with varying color properties.

\begin{figure}[!h]
	 \centerline{\includegraphics[width=0.8\textwidth]{src/augmentations/ColorJitter.png}}
	 {\caption{Examples of Color Jitter with the random changes in brightness, contrast, saturation, hue between 0 and 0.2} \label{fig:ColorJitter}}
\end{figure}

Similar but more sophisicated than \hyperlink[\ref{ChannelShuffles}]{Channel Shuffling}, Color Jitter can be advantageous for aerial images because it introduces random color variation to the images. For example, when detecting marine wildlife in aerial images, variations in the lighting and coloring of the water and surrounding environment can make it challenging to detect and differentiate the animals. By using Color Jitter to generate a diverse set of training images with random color perturbations, the machine learning model can learn to detect marine wildlife in a wider range of visual conditions. Additionally, standardizing the color distribution across images in the data-set can help to reduce the impact of color biases in the data-set, which can lead to improved accuracy and performance of the machine learning model, especially in classes with small samples. Figure \ref{fig:ColorJitter} illustrates the effect of the transformations.


\subsubsection{Solarize} \label{Solarize}
Solarize can be advantageous for aerial images, as it introduces a striking visual effect by inverting the colors of an image that exceed a certain threshold.

In the context of aerial images of marine wildlife, solarize can be useful to create a diverse set of training images for machine learning models. By inverting the colors of the images, solarize can create visually distinct images that can improve the model's ability to detect and recognize objects with varying color properties. Similar to Color Jitter \ref{ColorJitter}, Solarize helps to reduce the impact of color biases in the dataset. For example, if the majority of images of a particular marine species are predominantly of a particular color, the machine learning model may learn to associate that color with that class, leading to reduced accuracy and performance when detecting that species in images with different colors. By using solarize to invert the colors of these images, the model can be trained to recognize the species regardless of the color of the image.
\begin{figure}[!]
	 \centerline{\includegraphics[width=0.8\textwidth]{src/augmentations/Solarize.png}}
	 {\caption{Examples of Solarize with threshold of 128} \label{fig:cleah}}
\end{figure}

% \subsubsection{Normalize} \label{Normalize}

% Normalization is an image augmentation technique that involves scaling the pixel values of an image so that they fall within a specified range. This can be done by subtracting the mean pixel value of the image and dividing the result by the standard deviation of the pixel values. The concrete formula is:

% \begin{equation}
% img = \frac{(img - mean \cdot max\_pixel\_value)}{(std \cdot max\_pixel\_value)}
% \end{equation}

% As shown in in Figure \ref{fig:Normalize}, by scaling the pixel values to a standardized range, the machine learning model can better distinguish the features and patterns in the image that are relevant to the task, while reducing the impact of irrelevant variations in the image data, hence it can help to reduce the impact of variations in the lighting and color properties of the image.

% \begin{figure}[!h]
% 	 \centerline{\includegraphics[width=0.8\textwidth]{src/augmentations/Normalize.png}}
% 	 {\caption{Examples of Normalize Augmentation with the means respectively for three channels are 0.485, 0.456 and 0.406 and standard deviations are 0.229, 0.224 and 0.225.} \label{fig:Normalize}}
% \end{figure}

\subsubsection{ToGray} \label{ToGray}

One advantage of converting images into grayscale is that it can help to reduce the impact of color biases in the dataset. Color biases can occur when images of a particular class are predominantly of a particular color. By converting the images into grayscale, the model can be trained to recognize objects of a particular class based on their shapes and textures, rather than their color. However, the obvious disadvantage of converting images into grayscale is that it can result in the loss of important color information. Color Patterns may still be important for distinguishing different types of marine wildlife, and for identifying specific features of the environment that are important for object recognition. Figure \ref{fig:ToGray} illustrates our examples in grayscale.

\begin{figure}[!h]
	 \centerline{\includegraphics[width=0.8\textwidth]{src/augmentations/ToGray.png}}
	 {\caption{Examples of ToGray} \label{fig:ToGray}}
\end{figure}

\subsubsection{ToSepia} \label{ToSepia}

\begin{figure}
	 \centerline{\includegraphics[width=0.8\textwidth]{src/augmentations/ToSepia.png}}
	 {\caption{Examples of ToSepia} \label{fig:ToSepia}}
\end{figure}
Sepia is a type of image colorization technique that adds a brownish or yellowish tint to an image, giving it a vintage or aged look. It is often used in artistic or design applications to give images a nostalgic or classic feel.

While converting an image to grayscale  \ref{ToGray} involves discarding color information and representing the image in shades of gray, applying a sepia tint preserves the original color information and modifies it to create a vintage or aged effect. This makes the sepia technique a good option for applications such as edge detection, where the shape and structure of objects in the image are more important than their color, without the risk of significant loss of color information.

\subsubsection{Gaussian Noise} \label{GaussNoise}
\begin{figure}
	 \centerline{\includegraphics[width=0.8\textwidth]{src/augmentations/GaussNoise.png}}
	 {\caption{Examples of Gaussian Noise Injection with variance range for noise between 10 and 50} \label{fig:GaussNoise}}
\end{figure}

Gaussian noise injection is an image augmentation technique that adds random noise to the image by sampling from a normal distribution. In this technique, a random value is added to each pixel of the image, and the value of the added noise is determined by the standard deviation of the normal distribution, as illustrated in Figure \ref{fig:GaussNoise}.

The advantage of Gaussian noise injection is that it can create a diverse set of training images with varying levels of noise. This can help the machine learning model to better recognize objects in images with varying levels of noise, which can occur due to factors such as lighting, camera settings, or weather conditions. By training the model on a diverse set of noisy images, it can learn to better generalize to new, unseen data with varying levels of noise.

However, in the context of aerial images of marine wildlife, Gaussian noise injection may not be the most effective image augmentation technique. Aerial images of marine environments are often captured from a high altitude and are subject to atmospheric distortions, which can result in noise that is different from the type of noise added by Gaussian noise injection. Gaussian noise injection can negatively impact the visual quality of the image, which can make it more challenging to detect and recognize objects in the image.
\subsection{Cutout, GridDropout, Random Erasing}
In this section, we will introduce three augmentation techniques: Cutout, GridDropout, and Random Erasing. These techniques are frequently used to enhance the model's ability to concentrate on the important information in the image and increase its resistance to occlusion or missing data.
\subsubsection{Cutout/Coarse Dropout} \label{Cutout} 
Cutout is an image augmentation technique that involves randomly selecting a rectangular region of an image and setting the pixel values in that region to zero or a constant value. This creates a "hole" in the image, which can be used to encourage the model to focus on the remaining information in the image and improve its robustness to occlusion or missing data. The size, shape, and number of cutout regions can be varied to create a diverse set of training images and prevent overfitting.

Cutout augmentation can potentially be non label-preserving in aerial images because it involves removing parts of the image, which can also remove relevant objects or features that the model needs to learn, e.g a bird in the center of the image. This can cause inconsistencies between the original labels of the image and the labels assigned to the modified image, especially if the removed region contains the object of interest. For examples in Figure\ref{fig:Cutout}, one common goldeneye got cut out, if there were just one bird in the picture, it would be considered non-label preserving.

\begin{figure}[!h]
	 \centerline{\includegraphics[width=0.8\textwidth]{src/augmentations/Cutout.png}}
	 {\caption{Examples of Cutout with 8 boxes, each size 1/10 height of the image} \label{fig:Cutout}}
\end{figure}

\subsubsection{GridDropout}

GridDropout is an image augmentation technique that is based on the Cutout method \ref{Cutout} by the same author \cite{chen2020gridmask}. It involves dividing the image into a grid of equal-sized cells and randomly dropping out entire cells instead of rectangular patches. This creates a checkerboard pattern of dropped cells that encourages the model to learn features that are invariant to the exact location of objects in the image. The advantage of using GridDropout over Cutout is that it allows for larger and more complex holes to be created while also being less likely to remove entire objects from the image. In the case of our aerial image dataset, the objects of interest are often small in size, and therefore applying cutout augmentation may not be label-preserving due to the risk of removing important object features, as illustrated in Figure\ref{fig:GridDropout}.

\begin{figure}[!h]
	 \centerline{\includegraphics[width=0.8\textwidth]{src/augmentations/GridDropout.png}}
	 {\caption{Examples of GridDropout with 100 boxes} \label{fig:GridDropout}}
\end{figure}
\subsection{Mixed technique} 
This section will review shortly about the mixed approaches in image augmentation. Custom mixed augmentation and RandAugment will be introduced.
\subsubsection{RandAugment}\label{RandAugment}
RandAugment is a random image augmentation technique that applies a sequence of randomly selected image transformations to the training images \cite{cubuk2020randaugment}. The specific transformations and the sequence of their application are randomly generated for each image using a set of hyperparameters, resulting in a diverse set of training images with a wide range of visual variations. This can improve the performance of machine learning models by helping them to better generalize to new, unseen data, but careful tuning of the hyperparameters is necessary to avoid overfitting or underfitting the model. RandAugment is seen as a strong and useful way to increase data
for machine learning, and it is still being used today.

\begin{algorithm}[H]
\caption{RandAugment}\label{algo:randaugment}
% \begin{algorithmic}[1]
\Procedure{RandAugment}{$image, n, m$}
    \State $trans_list \gets$ a list of possible image transformations
    \State $trans_seq \gets$ a random sequence of $n$ transformations from $trans_list$
\For{$transform \in trans_seq$}
    \State $mag \gets$ a random number between 0 and $m$
    \State Apply the transformation $transform$ to the $image$ with magnitude $mag$
\EndFor
\State \textbf{return} the augmented image
\EndProcedure
% \end{algorithmic}
\end{algorithm}

This pseudo code \ref{algo:randaugment} of RandAugment  takes an input image, the number of transformations to apply ($n$), and the magnitude of the transformations ($m$) as parameters. The procedure first defines a list of possible image transformations and then generates a random sequence of $n$ transformations from this list. For each transformation in the sequence, the procedure applies the transformation to the input image with a randomly selected magnitude between 0 and $m$. Finally, the procedure returns the augmented image.

These specifics will depend on the application and dataset and may need to be adjusted accordingly. In our experiment, RandAugment takes transformation of \hyperref[Rotation]{Rotation}, \hyperref[Flipping]{Flipping}, \hyperref[MotionBlur]{Motion Blur}, \hyperref[Sharpen]{Sharpen}, \hyperref[CenterCrop]{Center Crop 0.4}, \hyperref[GaussNoise]{Gaussian Noise Injection}, \hyperref[ChannelShuffles]{Channel Shuffling}, \hyperref[{Cutout}]{Cutout}, and \hyperref[ColorJitter]{Color Jittering}. The samples of augmented images are illustrated in figure \ref{fig:RandAugment}

\begin{figure} [!h]
	 \centerline{\includegraphics[width=0.8\textwidth]{src/augmentations/RandAug.png}}
	 {\caption{Examples of RandAugment} \label{fig:RandAugment}}
\end{figure}
% \subsection{GAN-based augmentation}

% GAN-generated methods can help in augmentation by generating synthetic data that can be used to supplement the original training data, especially in cases where the training data is limited or imbalanced. The GAN model \ref{bg:gan} can learn the underlying distribution of the training data and generate new data samples that are similar to the real data but have different features that can help the model generalize better. These synthetic samples can be used to augment the training set and improve the model's performance. Additionally, GAN-generated methods can also be used to improve performance on edge cases, which are often underrepresented in the training data. By generating synthetic data for these edge cases, the model can be trained to better handle them, leading to improved overall performance.
% '
% TODO: Explain how GAN learn and helps.